\newcommand{\cauc@fakebold@size}{4}
\let\FakeBoldSize\cauc@fakebold@size
\newcommand{\cauc@fangsong@boldsize}{2.167}

\newif\ifcauc@localfonts@available
\cauc@localfonts@availabletrue

\def\cauc@localfonts@detect{%
  \IfFileExists{fonts/simsun.ttf}{}{\cauc@localfonts@availablefalse}%
  \IfFileExists{fonts/simhei.ttf}{}{\cauc@localfonts@availablefalse}%
  \IfFileExists{fonts/kaiti.ttf}{}{\cauc@localfonts@availablefalse}%
  \IfFileExists{fonts/lisu.ttf}{}{\cauc@localfonts@availablefalse}%
}

\ifx\cauc@localfonts\@empty
  \cauc@localfonts@detect
\else
  \def\cauc@@temp{auto}%
  \ifx\cauc@localfonts\cauc@@temp
    \cauc@localfonts@detect
  \else
    \def\cauc@@temp{true}%
    \ifx\cauc@localfonts\cauc@@temp
      \cauc@localfonts@detect
      \ifcauc@localfonts@available\else
        \ClassError{cauc_thesis}{%
          localfonts=true but font files not found in fonts/ directory.%
          \MessageBreak
          Please ensure fonts/simsun.ttf, fonts/simhei.ttf, fonts/kaiti.ttf, fonts/lisu.ttf exist.%
          \MessageBreak
          Falling back to system fonts.%
        }{%
          Place the required font files in the fonts/ directory, or set localfonts=false to use system fonts.%
        }%
      \fi
    \else
      \def\cauc@@temp{false}%
      \ifx\cauc@localfonts\cauc@@temp
        \cauc@localfonts@availablefalse
      \else
        \cauc@warning{Unknown localfonts value '\cauc@localfonts', expected: auto, true, false. Using auto.}%
        \cauc@localfonts@detect
      \fi
    \fi
  \fi
\fi

\ifcauc@localfonts@available
  \IfFileExists{fonts/times.ttf}{}{\cauc@localfonts@availablefalse}
  \IfFileExists{fonts/timesbd.ttf}{}{\cauc@localfonts@availablefalse}
  \IfFileExists{fonts/timesi.ttf}{}{\cauc@localfonts@availablefalse}
  \IfFileExists{fonts/timesbi.ttf}{}{\cauc@localfonts@availablefalse}
\fi

\ifcauc@localfonts@available
  % 宋体用于正文，黑体用于标题，楷体用于摘要关键词，隶书用于封面装饰
  \setCJKfamilyfont{songti}{simsun.ttf}[Path=fonts/]
  \setCJKfamilyfont{heiti}{simhei.ttf}[Path=fonts/]
  \setCJKfamilyfont{kaiti}{kaiti.ttf}[Path=fonts/]
  \setCJKfamilyfont{lishu}{lisu.ttf}[Path=fonts/]

  \setCJKfamilyfont{songtibold}{simsun.ttf}[Path=fonts/, AutoFakeBold={\cauc@fakebold@size}]
  \setCJKfamilyfont{heitibold}{simhei.ttf}[Path=fonts/, AutoFakeBold={\cauc@fakebold@size}]
  \setCJKfamilyfont{lishubold}{lisu.ttf}[Path=fonts/, AutoFakeBold={\cauc@fakebold@size}]
  \setCJKfamilyfont{kaitibold}{kaiti.ttf}[Path=fonts/, AutoFakeBold={\cauc@fakebold@size}]

  \IfFileExists{fonts/simfang.ttf}{%
    \setCJKfamilyfont{fangsong}{simfang.ttf}[Path=fonts/, AutoFakeBold={\cauc@fangsong@boldsize}]%
    \setCJKmonofont{simfang.ttf}[Path=fonts/, AutoFakeBold={\cauc@fangsong@boldsize}]%
  }{%
    \setCJKfamilyfont{fangsong}{simsun.ttf}[Path=fonts/, AutoFakeBold={\cauc@fangsong@boldsize}]%
    \setCJKmonofont{simsun.ttf}[Path=fonts/, AutoFakeBold={\cauc@fangsong@boldsize}]%
    \cauc@warning{FangSong font not found, falling back to SimSun for fangsong/monofont}%
  }

  % 宋体无斜体变体，用楷体替代宋体斜体，避免Font shape undefined警告
  \setCJKmainfont[Path=fonts/, ItalicFont=kaiti.ttf]{simsun.ttf}
  % fontset=none下需显式设置CJK无衬线字体族，黑体作为CJK sans-serif默认字体
  \setCJKsansfont{simhei.ttf}[Path=fonts/]
  \setmainfont{times.ttf}[Path=fonts/, BoldFont=timesbd.ttf, ItalicFont=timesi.ttf, BoldItalicFont=timesbi.ttf]
\else
  % 优先加载本地fonts/目录字体以确保跨平台一致性，不存在时回退到系统字体
  \cauc@warning{Local font files not found in fonts/ directory. Falling back to system fonts.}

  \setCJKfamilyfont{songti}{SimSun}
  \setCJKfamilyfont{heiti}{SimHei}
  \setCJKfamilyfont{kaiti}{KaiTi}
  \setCJKfamilyfont{lishu}{LiSu}

  \setCJKfamilyfont{songtibold}{SimSun}[AutoFakeBold={\cauc@fakebold@size}]
  \setCJKfamilyfont{heitibold}{SimHei}[AutoFakeBold={\cauc@fakebold@size}]
  \setCJKfamilyfont{lishubold}{LiSu}[AutoFakeBold={\cauc@fakebold@size}]
  \setCJKfamilyfont{kaitibold}{KaiTi}[AutoFakeBold={\cauc@fakebold@size}]

  \IfFontExistsTF{FangSong}{%
    \setCJKfamilyfont{fangsong}{FangSong}[AutoFakeBold={\cauc@fangsong@boldsize}]%
    \setCJKmonofont{FangSong}[AutoFakeBold={\cauc@fangsong@boldsize}]%
  }{%
    \IfFontExistsTF{STFangSong}{%
      \setCJKfamilyfont{fangsong}{STFangSong}[AutoFakeBold={\cauc@fangsong@boldsize}]%
      \setCJKmonofont{STFangSong}[AutoFakeBold={\cauc@fangsong@boldsize}]%
    }{%
      \setCJKfamilyfont{fangsong}{SimSun}[AutoFakeBold={\cauc@fangsong@boldsize}]%
      \setCJKmonofont{SimSun}[AutoFakeBold={\cauc@fangsong@boldsize}]%
      \cauc@warning{FangSong font not found, falling back to SimSun for fangsong/monofont}%
    }%
  }

  % 宋体无斜体变体，用楷体替代宋体斜体，避免Font shape undefined警告
  \setCJKmainfont{SimSun}[ItalicFont=KaiTi]
  % fontset=none下需显式设置CJK无衬线字体族，黑体作为CJK sans-serif默认字体
  \setCJKsansfont{SimHei}
  \setmainfont{Times New Roman}
\fi

\newcommand{\heitibold}{\CJKfamily{heitibold}}
\newcommand{\songtibold}{\CJKfamily{songtibold}}
\newcommand{\lishubold}{\CJKfamily{lishubold}}
\newcommand{\kaitibold}{\CJKfamily{kaitibold}}
\let\fakehei\heitibold
\let\fakesong\songtibold
\let\fakeli\lishubold
\let\fakekai\kaitibold
% fontset=none下ctex不预定义字体命令，使用\DeclareDocumentCommand兼容"已存在"和"不存在"两种情况
\DeclareDocumentCommand{\songti}{}{\CJKfamily{songti}}
\DeclareDocumentCommand{\heiti}{}{\CJKfamily{heiti}}
\newcommand{\kaiti}{\CJKfamily{kaiti}}
\DeclareDocumentCommand{\lishu}{}{\CJKfamily{lishu}}
\DeclareDocumentCommand{\fangsong}{}{\CJKfamily{fangsong}}

\newcommand{\timu}{\fontsize{16pt}{\baselineskip}\selectfont\CJKfamily{heiti}}
\newcommand{\zhaiyao}{\fontsize{16pt}{\baselineskip}\selectfont\CJKfamily{lishu}}
